% Copyright (c) 2026 Eugene Joseph M. Ragasa.
% SPDX-License-Identifier: Apache-2.0
% Author: Eugene Joseph M. Ragasa
% Formal proof completed: 2026-10-20
\documentclass[11pt]{article}
\usepackage[T1]{fontenc}
\usepackage{amsmath,amssymb}
\usepackage[margin=1in]{geometry}
\usepackage{xurl}
\usepackage[hidelinks]{hyperref}
\title{Proof Dossier: Lifting a Test-Space Formal-Adjoint Identity}
\author{Eugene Joseph M. Ragasa}
\date{Formal proof completed: 2026-10-20}
\begin{document}
\maketitle

\section{Identity and status}
ID \texttt{QM-OPERATOR-THM-002}; status \texttt{proved; no admissions}.
Owning module: \path{ProjectKoios/QuantumMechanics/Operator/TestDomain.lean}.
Lean declaration:
\par\noindent{\scriptsize\path{ProjectKoios.QuantumMechanics.isFormalAdjoint_testSpaceOperator}}.

\section{Exact mathematical statement}
Let $i:\Phi\to\mathcal H$ be the injective dense inclusion of a
\texttt{RiggedHilbertSpaceCore}, and let $A,B:\Phi\to\Phi$ be complex linear maps. Suppose
\[
  \langle i(A\phi),i\psi\rangle
  =\langle i\phi,i(B\psi)\rangle
  \qquad\text{for every }\phi,\psi\in\Phi,
\]
using Mathlib's convention that the inner product is conjugate-linear in its first argument. Then
the induced partial Hilbert-space operators $T_A$ and $T_B$, both with domain $i(\Phi)$, satisfy
Mathlib's formal-adjoint relation
\[
  \langle T_A x,y\rangle=\langle x,T_B y\rangle
  \qquad(x\in\mathcal D(T_A),\ y\in\mathcal D(T_B)).
\]

\section{Source, assumptions, and statement delta}
This is source-neutral functional analysis. The displayed test-space identity is an explicit
hypothesis, not a repository axiom and not a conclusion derived from the name of an operator. The
theorem does not assert that $A$ or $B$ is continuous, symmetric, closed, closable, mutually
adjoint as maximal operators, or self-adjoint.

\section{Prior art and contribution boundary}
Passing an inner-product identity through an injective domain parametrization is standard. The
repository contribution is the typed bridge from test-space endomorphisms to Mathlib's
\path{LinearPMap.IsFormalAdjoint} predicate. No novelty or textbook-fidelity claim is made.

\section{Proof schematic}
For arbitrary domain vectors $x$ and $y$:
\[
 x=i\phi,\quad y=i\psi
 \Longrightarrow
 T_Ax=i(A\phi),\quad T_By=i(B\psi)
 \Longrightarrow
 \langle T_Ax,y\rangle=\langle x,T_By\rangle.
\]
Injectivity of $i$ makes the representatives selected by the range equivalence unique.

\section{Human-readable proof}
Choose the unique test vectors $\phi$ and $\psi$ corresponding to the domain elements $x$ and $y$
through the linear equivalence between $\Phi$ and $\operatorname{ran} i$. By construction,
$x=i\phi$ and $y=i\psi$. The induced-action identity gives $T_Ax=i(A\phi)$ and
$T_By=i(B\psi)$. Substitution reduces the desired formal-adjoint equality exactly to the supplied
hypothesis for $\phi$ and $\psi$.

\section{Lean realization and library dependencies}
The proof defines $\phi=(\texttt{testDomainEquiv }R)^{-1}x$ and similarly for $\psi$, uses the
left/right inverse laws of the linear equivalence, rewrites both partial-operator actions with
\path{testSpaceOperator_apply_includedTestVector}, and closes with
\path{IsTestSpaceFormalAdjoint}. The formal-adjoint target is Mathlib's existing partial-operator
predicate.

\section{Trusted base, automation, and toolchain}
The declaration is kernel checked with Lean \texttt{v4.35.0} and Mathlib revision
\path{56a005855ea17e0b0958f73ec801ae140b0a1679}. The recorded \texttt{\#print axioms} result is
\texttt{propext}, \texttt{Classical.choice}, and \texttt{Quot.sound}. There are no admissions or
repository-owned scientific axioms. AI assistance contributed to proof discovery, Lean
implementation, and exposition; no independent expert review is recorded.

\section{Evidence and applicability boundary}
The premise must be established separately for every concrete operator pair. Numerical agreement
cannot silently replace it. The result supplies no maximal-adjoint equality, self-adjointness,
observable interpretation, spectral theorem, generalized eigenvector, or physical applicability.

\section{Unresolved questions and revision history}
Concrete invariant test spaces, concrete operator pairs, and conditions upgrading formal-adjoint
relations remain open. Initial proof and dossier completed 2026-10-20.
\end{document}
